Having established the requirements of FreelanceChain in Chapter 2, this chapter presents the technologies and methodologies selected for the implementation. For each technology, the rationale for its selection over principal alternatives is discussed. Section \ref{section:3.1} covers the Solana blockchain and its consensus mechanism. Section \ref{section:3.2} presents the Anchor smart contract framework. Section \ref{section:3.3} describes the Rust programming language. Section \ref{section:3.4} covers the Next.js and React frontend stack. Section \ref{section:3.5} explains the face recognition technology used for KYC. Section \ref{section:3.6} describes the Groq-based AI risk assessment and conversational advisor integration. Section \ref{section:3.7} covers the \gls{ipfs} storage layer.

\section{Solana Blockchain}
\label{section:3.1}

\subsection{Overview and Consensus Mechanism}

Solana is a public, permissionless Layer-1 blockchain designed for high-performance decentralized applications \cite{yakovenko2018solana}. Its principal innovation is the Proof-of-History (PoH) mechanism, which establishes a cryptographic clock that orders events before they are submitted to the consensus process. By pre-determining the sequence of events, Solana's validators can process transactions in parallel without communicating about event ordering, dramatically increasing throughput.

In combination with Tower BFT (a PoH-optimized version of Practical Byzantine Fault Tolerance), Solana achieves a theoretical throughput of over 65,000 \gls{tps} with average block confirmation times of approximately 400 milliseconds. The practical throughput on the live network consistently exceeds 2,000 \gls{tps} under load, compared to Ethereum's practical limit of approximately 15--30 \gls{tps} \cite{solana2021performance}.

\subsection{Account Model}

Unlike Ethereum's account-centric model where contract code and state are stored together, Solana separates code from data. Programs (smart contracts) are stateless executable objects stored in program accounts; all state is stored in separate data accounts that programs can read and write. This separation enables parallel execution of transactions that access different accounts.

\textbf{Program Derived Addresses (\gls{pda}s)} are deterministic account addresses generated by hashing a set of seeds together with the program ID. PDAs have no corresponding private key, meaning they can only be signed by the program itself during instruction execution. This property is exploited extensively in FreelanceChain: every job, bid, escrow account, milestone, KYC record, reputation SBT, and governance proposal is stored in a distinct PDA derived from meaningful seeds (e.g., \texttt{[b"job", client\_pubkey]}).

\subsection{Solana vs. Ethereum: Rationale for Selection}

The primary alternative to Solana is Ethereum, which has the largest smart contract developer ecosystem and the most mature tooling. However, three factors favor Solana for FreelanceChain:

\begin{enumerate}
    \item \textbf{Transaction cost:} Solana's average transaction fee is below \$0.001, compared to Ethereum's fees that regularly exceed \$5--\$50 on mainnet. For a freelance platform where users may submit dozens of transactions per engagement (bids, milestone fundings, approvals), high fees would make the platform economically unviable for small-scale work.
    \item \textbf{Throughput:} Solana's throughput allows the platform to scale to millions of concurrent users without congestion, whereas Ethereum's 15--30 \gls{tps} limit would create bottlenecks at scale.
    \item \textbf{User experience:} Solana's sub-second confirmation times enable a user experience comparable to traditional web applications, whereas Ethereum's 15-second block times introduce perceptible delays.
\end{enumerate}

Polygon and Avalanche are Ethereum-compatible chains that offer lower fees and higher throughput, but they sacrifice some degree of decentralization and have smaller ecosystem tooling for the specific libraries (Anchor, SPL) used in FreelanceChain.

\section{Anchor Smart Contract Framework}
\label{section:3.2}

\subsection{Overview}

Anchor \cite{anchor2021} is a development framework for Solana programs written in Rust. It provides a set of macros and abstractions that handle the boilerplate required for Solana program development: account validation, deserialization/serialization of account data, error handling, and \gls{idl} generation.

A key feature of Anchor is its account constraint system. By annotating accounts with constraints such as \texttt{\#[account(init, payer = user, space = 8 + JobAccount::INIT\_SPACE)]} or \texttt{\#[account(mut, has\_one = authority)]}, the framework generates validation code that checks these constraints before any instruction logic executes. This dramatically reduces the risk of common smart contract vulnerabilities such as missing ownership checks, account substitution attacks, and integer overflow.

\subsection{Program Derived Addresses in Anchor}

Anchor's \texttt{\#[account(seeds = [...], bump)]} constraint automates PDA derivation and bump seed storage. This simplifies the process of creating and validating PDAs, and the \texttt{init\_if\_needed} feature allows an account to be initialized on first use, reducing the number of transactions required for onboarding.

\subsection{Interface Definition Language}

Anchor automatically generates an \gls{idl} file from the Rust source code annotations. The IDL is a JSON schema that describes all program instructions (name, accounts, arguments, types) and account structures. This IDL is used by the TypeScript \gls{sdk} (\texttt{@coral-xyz/anchor}) on the frontend to construct and send transactions without manually constructing the raw binary instruction format.

\subsection{Anchor vs. Native Solana Programs: Rationale for Selection}

Writing native Solana programs in raw Rust requires manually implementing all account validation logic, deserialization, and PDA arithmetic. This increases the probability of security vulnerabilities and dramatically extends development time. Anchor's constraint system provides compile-time guarantees that reduce the attack surface and allows the developer to focus on business logic rather than low-level account handling. For the eighteen modules required by FreelanceChain, Anchor was the clear choice.

\section{Rust Programming Language}
\label{section:3.3}

Rust is a systems programming language designed for performance and safety \cite{rust2023}. Its ownership model enforces memory safety at compile time without a garbage collector, eliminating entire classes of bugs (use-after-free, null dereference, data races) that commonly plague smart contracts written in other languages. Solana programs must be compiled to Berkeley Packet Filter (\gls{sdk}) bytecode, and Rust is the only officially supported language for this compilation target.

FreelanceChain's eighteen Anchor modules are each implemented as a Rust module within a single Cargo workspace. The program binary was size-optimized to approximately 1.6 MB using Cargo's \texttt{opt-level = "z"} setting and by removing unused dependency features (specifically the \texttt{mpl-token-metadata} and \texttt{anchor-spl/metadata} features), reducing the program data account size from over 2 MB.

\section{Next.js and React Frontend Stack}
\label{section:3.4}

\subsection{React and Component Architecture}

React \cite{react2023} is a declarative JavaScript library for building \gls{ui} components. Its unidirectional data flow and component model provide a predictable state management pattern that scales well to complex applications. FreelanceChain's frontend comprises over 50 React components organized into pages, layout components, feature-specific components, and reusable primitive components.

State management is handled by TanStack Query (React Query) \cite{reactquery2023}, which provides a robust caching, synchronization, and re-fetching layer for asynchronous data from both the Solana blockchain and the Next.js \gls{api} routes. This eliminates the need for a dedicated global state manager (such as Redux) and ensures that blockchain state is always reflected accurately in the \gls{ui}.

\subsection{Next.js}

Next.js \cite{nextjs2023} is a React meta-framework that provides server-side rendering, static site generation, file-based routing, and built-in \gls{api} route support. FreelanceChain uses Next.js for:

\begin{enumerate}
    \item \textbf{File-based routing:} Pages such as \texttt{/jobs/[id]}, \texttt{/profile/[address]}, and \texttt{/identity} are automatically routed based on the filesystem structure.
    \item \textbf{API routes:} Serverless functions at \texttt{/api/ai/risk-assessment}, \texttt{/api/cv-proxy}, and \texttt{/api/ai/risk-chat} handle AI integration and \gls{ipfs} proxying server-side.
    \item \textbf{Server-side rendering:} Job detail pages can be pre-rendered with public blockchain data for improved SEO and initial load performance.
\end{enumerate}

\subsection{Material UI}

Material UI (MUI) \cite{materialui2023} provides a comprehensive library of pre-built React components adhering to Google's Material Design specification. For FreelanceChain, MUI's component library accelerates frontend development by providing production-ready implementations of tables, modals, drawers, tabs, tooltips, and form controls. The MUI theming system allows a custom crypto-terminal aesthetic (dark background, \texttt{\#00ffc3} teal accent, Orbitron display font) to be applied globally without modifying individual components.

\subsection{Framer Motion}

Framer Motion \cite{framermotion2023} provides a declarative animation API for React. FreelanceChain uses it for page transitions, hover effects on interactive elements, and ambient animations on the homepage that enhance the premium crypto-platform aesthetic without impacting performance.

\subsection{Next.js vs. Create React App: Rationale for Selection}

Create React App (CRA) produces client-side-only applications without server-side rendering or API routes. FreelanceChain requires both --- server-side rendering for public pages and API routes for AI integration --- making Next.js the correct choice. Vite was considered as an alternative for its superior development-server performance, but the lack of built-in API routes would have required a separate Express server, increasing deployment complexity.

\section{Face Recognition for KYC}
\label{section:3.5}

\subsection{face-api.js and Model Architecture}

FreelanceChain integrates the \texttt{@vladmandic/face-api} library \cite{vladmandic2021faceapi}, a maintained fork of the original \texttt{face-api.js}, which provides face detection and recognition capabilities running entirely in the browser via TensorFlow.js. Three neural network models are used:

\begin{enumerate}
    \item \textbf{TinyFaceDetector:} A lightweight convolutional neural network (CNN) that detects face bounding boxes in an image. Its small size (~190 KB) enables rapid inference even on consumer hardware.
    \item \textbf{FaceLandmark68TinyNet:} Identifies 68 facial landmark points (eyes, nose, mouth, jawline) that are used to align the face before descriptor computation.
    \item \textbf{FaceRecognitionNet:} A ResNet-34--based architecture that maps an aligned face image to a 128-dimensional feature vector (face descriptor). This descriptor encodes the unique geometric characteristics of a face in a compact numerical representation.
\end{enumerate}

\subsection{Verification Algorithm}

Given a face descriptor extracted from an identity document image ($\mathbf{d}_{\text{doc}}$) and a descriptor from a selfie ($\mathbf{d}_{\text{selfie}}$), the Euclidean distance is computed as:

\begin{equation}
    \delta = \left\| \mathbf{d}_{\text{doc}} - \mathbf{d}_{\text{selfie}} \right\|_2 = \sqrt{\sum_{i=1}^{128} (d_{\text{doc},i} - d_{\text{selfie},i})^2}
\end{equation}

If $\delta < 0.50$, the two faces are deemed to belong to the same person and the verification is approved. The threshold of 0.50 was selected to balance false-rejection rate (too strict) against false-acceptance rate (too permissive). The face-api.js documentation recommends thresholds between 0.45 and 0.60; 0.50 was adopted as the midpoint, accommodating natural variation in lighting, camera angle, and aging between the identity document photo and the live selfie.

On-chain, the distance is stored as basis points (\texttt{face\_distance\_bp} = $\delta \times 10{,}000$ as a \texttt{u32}), avoiding floating-point arithmetic in the Rust smart contract.

\subsection{Privacy Design}

All face recognition computation occurs in the user's browser. No image and no feature vector is transmitted to any server. The only data written to the blockchain is:
\begin{itemize}
    \item The identity document type (national ID, passport, or driver's licence).
    \item The face distance in basis points.
    \item A boolean matched flag.
    \item A status enum (Verified / Rejected).
\end{itemize}

This design is privacy-preserving by construction: even if the blockchain were fully public (which it is), no biometric information is recoverable from the stored values.

\subsection{face-api.js vs. Server-Side Facial Recognition: Rationale}

Server-side face recognition (e.g., AWS Rekognition, Azure Face) would provide higher accuracy and more robust illumination correction. However, it would require transmitting face images to a third-party server, creating a centralized biometric data store that is antithetical to the decentralized privacy principles of FreelanceChain. Browser-side inference with face-api.js achieves adequate accuracy (reported equal error rate of 97.3\% on LFW benchmark for the ResNet-34 model) while eliminating all biometric data custody.

\section{AI Risk Assessment: Groq API and Llama 3.3 70B}
\label{section:3.6}

\subsection{Overview of the AI Stack}

FreelanceChain integrates open-source large language model inference via the Groq \gls{api} using the \texttt{llama-3.3-70b-versatile} model. Groq is a specialized AI inference provider that runs open-source models (including Meta's LLaMA family) on custom-designed Language Processing Units (LPUs), achieving substantially lower latency than GPU-based inference. The integration consists of two complementary serverless API routes implemented in Next.js:

\begin{enumerate}
    \item \texttt{/api/ai/risk-assessment} --- performs structured multi-dimensional analysis of a freelancer's CV against the job requirements, returning a composite risk score and classification.
    \item \texttt{/api/ai/risk-chat} --- provides a conversational follow-up advisor that answers client questions about the risk assessment using the full prior assessment as context.
\end{enumerate}

\subsection{Four-Score Risk Assessment Model}

The \texttt{risk-assessment} route implements a four-dimensional scoring model. When a client submits a freelancer's CV text and the job description, the API route sends a structured prompt to \texttt{llama-3.3-70b-versatile} instructing it to independently evaluate four dimensions:

\begin{enumerate}
    \item \textbf{matchScore (0--100):} Degree of alignment between the freelancer's declared skills, experience, and portfolio and the specific requirements of the job posting.
    \item \textbf{authenticityScore (0--100):} Assessment of CV credibility indicators --- consistency of employment timelines, specificity of project descriptions, plausibility of claimed achievements, and absence of generic filler language.
    \item \textbf{budgetScore (0--100):} Evaluation of whether the freelancer's proposed bid amount is realistic and reasonable relative to the scope, complexity, and market rate implied by the job description.
    \item \textbf{timelineScore (0--100):} Evaluation of whether the freelancer's proposed delivery timeline is achievable given the scope of work and the skills indicated in the CV.
\end{enumerate}

These four scores are combined into a composite score using a weighted formula that prioritizes skill alignment and CV authenticity over financial parameters:

\begin{equation}
    S_{\text{composite}} = 0.35 \cdot S_{\text{match}} + 0.35 \cdot S_{\text{auth}} + 0.15 \cdot S_{\text{budget}} + 0.15 \cdot S_{\text{timeline}}
\end{equation}

The risk score is defined as the complement:

\begin{equation}
    S_{\text{risk}} = 100 - S_{\text{composite}}
\end{equation}

Risk classification thresholds are applied to $S_{\text{risk}}$:
\begin{itemize}
    \item $S_{\text{risk}} \leq 35$: \texttt{LOW} risk --- candidate is a strong match.
    \item $35 < S_{\text{risk}} \leq 60$: \texttt{MEDIUM} risk --- candidate has relevant skills but notable concerns.
    \item $S_{\text{risk}} > 60$: \texttt{HIGH} risk --- significant mismatch or authenticity concerns.
\end{itemize}

The model's JSON response is parsed using a primary structured extractor; if the model response is not strictly valid JSON (a failure mode observed in $\approx$5\% of calls), a regex-based fallback extracts the four score fields from the raw response text, ensuring robustness.

\subsection{PDF Text Extraction Pipeline}

To reduce friction for clients, the \texttt{risk-assessment} route also supports PDF uploads. A two-stage extraction pipeline processes uploaded PDF buffers:

\begin{enumerate}
    \item \textbf{Primary (unpdf / pdfjs-dist):} The \texttt{unpdf} library, built on Mozilla's PDF.js engine, extracts text from the PDF's content stream with correct Unicode handling and layout-aware paragraph ordering.
    \item \textbf{Fallback (zlib + regex):} If the primary extractor fails (e.g., scanned PDFs or non-standard PDF encodings), the raw PDF buffer is decompressed using Node.js's \texttt{zlib} module, and a regex pattern extracts readable text spans from the decompressed content stream. This provides a best-effort extraction for edge-case documents.
\end{enumerate}

The extracted text is then passed to the LLM as the CV input, enabling a fully server-side pipeline where the client need only upload a PDF file.

\subsection{Conversational AI Advisor}

The \texttt{risk-chat} route extends the risk assessment with an interactive conversational advisor. After the initial risk assessment is displayed, the client can ask follow-up questions in a chat interface (e.g., ``What specific skills are missing?'' or ``Is the proposed budget fair for a senior developer?''). The API route:

\begin{enumerate}
    \item Constructs a system prompt that embeds the full prior assessment context: all four scores, the key findings, the risk classification, and the recommendation.
    \item Appends the conversation history (prior user questions and model replies).
    \item Sends the complete context to \texttt{llama-3.3-70b-versatile} with instructions to respond as a professional hiring advisor in 2--5 concise sentences.
\end{enumerate}

This design keeps the conversational AI grounded in the specific assessment data, preventing generic or hallucinated responses about the candidate.

\subsection{Groq / Llama 3.3 vs. Alternatives: Rationale}

\texttt{llama-3.3-70b-versatile} via the Groq API was selected over OpenAI GPT-4o, Google Gemini 1.5 Flash, and the Anthropic Claude API on four grounds:

\begin{enumerate}
    \item \textbf{Inference speed:} Groq's LPU hardware achieves token generation rates exceeding 500 tokens/second for LLaMA 70B models, yielding median response times below 2 seconds for the prompt sizes used. This comfortably meets the 10-second \gls{ux} requirement even for lengthy CV inputs.
    \item \textbf{Cost:} The Groq API offers a free tier sufficient for development and demonstration workloads, with no credit card required. Production pricing is competitive with other providers at comparable quality levels.
    \item \textbf{Model capability:} LLaMA 3.3 70B achieves state-of-the-art performance on instruction-following and structured output benchmarks among open-source models, producing well-formed JSON reliably with the structured prompt design used in this work.
    \item \textbf{Open-source alignment:} Using an open-source model (Meta LLaMA) aligns with the decentralization philosophy of FreelanceChain; the model weights can be self-hosted to eliminate the remaining off-chain AI dependency in future work.
\end{enumerate}

\section{IPFS for Decentralized Storage}
\label{section:3.7}

\gls{ipfs} \cite{ipfs2015} is a peer-to-peer distributed file system that identifies content by its cryptographic hash (Content ID, or CID) rather than by location. This content-addressing property means that a file's CID is stable regardless of which node stores it, and the content is verifiable against the CID at retrieval time.

FreelanceChain uses IPFS for:
\begin{itemize}
    \item Job description metadata (title, description, required skills, attachments).
    \item CV documents uploaded by freelancers with bids.
    \item Work submission deliverables.
    \item Dispute evidence packages.
    \item SBT metadata URIs for reputation tokens.
\end{itemize}

Only the CID (a short string) is stored on-chain, minimizing the on-chain data footprint and cost. The full document is retrieved client-side using an IPFS gateway at read time. This architecture cleanly separates mutable large-form content (IPFS) from immutable transactional records (Solana blockchain).

In summary, the technology choices for FreelanceChain form a coherent, mutually reinforcing stack: Solana's high throughput and low fees make frequent blockchain interactions economically viable; Anchor's constraint system provides security guarantees at the smart contract layer; Next.js enables a production-quality web frontend; face-api.js delivers privacy-preserving identity verification; the Groq API with LLaMA 3.3 70B adds AI-assisted candidate evaluation and conversational advisory; and IPFS provides decentralized document storage. Each choice was made after a comparative analysis of alternatives, with the selection criteria being technical fitness, economic viability, developer productivity, and alignment with the decentralization principles of the platform.
